\documentclass[a4,11pt]{aleph-notas}
% Se puede ver la documentación aquí:
% https://github.com/alephsub0/LaTeX_aleph-notas

% -- Paquetes adicionales
\usepackage{enumitem}
\usepackage{url}
\usepackage{array}
\usepackage{booktabs}
\usepackage{longtable}
\hypersetup{
    urlcolor=blue,
    linkcolor=blue,
}

% -- Datos
\institucion{Facultad de Ciencias Exactas, Naturales y Ambientales}
\carrera{Catálogo STEM}
\asignatura{Álgebra lineal}
\tema{Clase 02: Sistemas de ecuaciones lineales}
\autor{Andrés Merino}
\fecha{Periodo 2026-01}

\logouno[0.16\textwidth]{Logos/logoPUCE_04_ac}
\definecolor{colortext}{HTML}{1957d1}
\definecolor{colordef}{HTML}{1957d1}
\fuente{montserrat}

\begin{document}

\encabezado

%%%%%%%%%%%%%%%%%%%%%%%%%%%%%%%%%%%%%%%%
\section*{Resultado de Aprendizaje}
%%%%%%%%%%%%%%%%%%%%%%%%%%%%%%%%%%%%%%%%

%%%%%%%%%%%%%%%%%%%%%%%%%%%%%%%%%%%%%%%%
\subsection*{RdA de la asignatura:}
\begin{itemize}[leftmargin=*]
    \item \textbf{RdA 1:} Comprender los conceptos fundamentales del Álgebra Lineal, incluyendo el estudio de matrices, determinantes, sistemas de ecuaciones lineales y espacios vectoriales, destacando su importancia en el análisis y resolución de problemas matemáticos.
    \item \textbf{RdA 2:} Resolver problemas prácticos y teóricos relacionados con el Álgebra Lineal, utilizando técnicas de cálculo con matrices, determinantes y sistemas de ecuaciones, así como en el análisis de espacios vectoriales.
\end{itemize}

%%%%%%%%%%%%%%%%%%%%%%%%%%%%%%%%%%%%%%%%
\subsection*{RdA de la actividad:}
\begin{itemize}[leftmargin=*]
    \item Identificar la forma matricial de un sistema de ecuaciones lineales.
    \item Aplicar operaciones elementales por filas para obtener formas escalonadas.
    \item Determinar si un sistema es consistente, inconsistente o con infinitas soluciones.
\end{itemize}

%%%%%%%%%%%%%%%%%%%%%%%%%%%%%%%%%%%%%%%%
\section*{Introducción}
%%%%%%%%%%%%%%%%%%%%%%%%%%%%%%%%%%%%%%%%

%%%%%%%%%%%%%%%%%%%%%%%%%%%%%%%%%%%%%%%%
\paragraph{Control de lectura:}
Antes de iniciar la clase, se aplicará el control de lectura del siguiente enlace:
\href{https://gemini.google.com/share/def9847cb644}{Control de lectura - Clase 02}.

\paragraph{Pregunta inicial:}
¿Un sistema con el mismo número de ecuaciones que incógnitas siempre tiene solución? ¿Por qué a veces no?

%%%%%%%%%%%%%%%%%%%%%%%%%%%%%%%%%%%%%%%%
\section*{Desarrollo}
%%%%%%%%%%%%%%%%%%%%%%%%%%%%%%%%%%%%%%%%

%%%%%%%%%%%%%%%%%%%%%%%%%%%%%%%%%%%%%%%%
\subsection*{Actividad 1: Resolución de sistemas mediante eliminación de Gauss-Jordan}

La clase inicia con un problema contextual que conduce a un sistema de ecuaciones lineales. Luego, mediante clase magistral y práctica guiada, se trabaja la representación matricial, la matriz ampliada y el proceso de eliminación de Gauss-Jordan para analizar el tipo de solución.

\paragraph{¿Cómo lo haremos?}
\begin{itemize}[leftmargin=*]
    \item \textbf{Motivación:} se presenta una situación aplicada (por ejemplo, mezcla o balance) y se plantea el sistema lineal correspondiente.
    \item \textbf{Clase magistral:} se explican matriz de coeficientes, matriz ampliada, operaciones elementales por filas, forma escalonada y forma escalonada reducida. Se usa el resumen \href{https://fcena-puce.github.io/AlgebraLineal-05-N0048/01Resumen/Resumen02.pdf}{Resumen02.pdf}.
    \item \textbf{Resolución de ejercicios:} los estudiantes resolverán, de manera guiada, ejercicios de clasificación y resolución de sistemas lineales del documento \href{https://fcena-puce.github.io/AlgebraLineal-05-N0048/04Ejercicios/Ejercicios02.pdf}{Ejercicios02.pdf}.
    \item \textbf{Visualización de videos:} se recomendará el estudio con los videos:  
    \begin{itemize}
        \item \href{https://www.youtube.com/watch?v=IF7WV5VVci4}{Sistemas de dos Ecuaciones}
        \item \href{https://www.youtube.com/watch?v=dFmGzr1j6eY}{Solución de sistemas de 3x3 método de Gauss Jordan}
    \end{itemize}  
    \item \textbf{Implementación computacional:} se realiza mediante cuadernos de Jupyter usando el documento \href{https://fcena-puce.github.io/AlgebraLineal-05-N0048/02ResumenPython/ResumenPy02.pdf}{ResumenPy02.pdf}.
\end{itemize}

\paragraph{Verificación de aprendizaje:}
\begin{itemize}[leftmargin=*]
    \item ¿Qué diferencia existe entre la matriz de coeficientes y la matriz ampliada?
    \item ¿Cómo se interpreta una fila de la forma $[0\ 0\ \cdots\ 0\ |\ b]$ con $b\neq 0$?
    \item ¿Qué criterio permite decidir si un sistema tiene solución única o infinitas soluciones?
\end{itemize}

%%%%%%%%%%%%%%%%%%%%%%%%%%%%%%%%%%%%%%%%
\section*{Cierre}
%%%%%%%%%%%%%%%%%%%%%%%%%%%%%%%%%%%%%%%%

\paragraph{Tarea:}  
Resolver del libro \href{https://catalogobiblioteca.puce.edu.ec/cgi-bin/koha/opac-detail.pl?biblionumber=86081}{Fundamentos de álgebra lineal de Larson}, sección 1.2, los ejercicios: 11, 13, 25, 29, 30, 33, 35, 37, 49, 50.

\paragraph{Pregunta de investigación:}
\begin{enumerate}[leftmargin=*]
    \item ¿Qué papel juegan los sistemas lineales en el entrenamiento de modelos de inteligencia artificial?
    \item ¿Qué significa que una matriz tenga rango completo? ¿Y qué consecuencias tiene para la resolución de sistemas?
\end{enumerate}

\paragraph{Para la próxima clase:}  
Leer la sección 2.2 del libro \href{https://catalogobiblioteca.puce.edu.ec/cgi-bin/koha/opac-detail.pl?biblionumber=283260}{Álgebra lineal con aplicaciones y Python de Aranda} y la sección 2.3 del libro \href{https://catalogobiblioteca.puce.edu.ec/cgi-bin/koha/opac-detail.pl?biblionumber=86081}{Fundamentos de álgebra lineal de Larson}.

\end{document}
